\section{Introduction}
\label{sec:intro}

The Liechtenstein-Katsnelson-Antropov-Gubanov (LKAG) construction maps infinitesimal electronic spin rotations onto parameters of an effective spin Hamiltonian using the magnetic force theorem. Its nonrelativistic and relativistic forms are now standard tools for exchange, anisotropic exchange, and Dzyaloshinskii-Moriya interactions~\cite{Liechtenstein1987,Udvardi2003,EbertMankovsky2009,Szilva2023}. Localized-orbital implementations are attractive because the magnetic degrees of freedom are already represented in a compact basis, but the definition of a local operator and the choice of the rotated magnetic subspace require care~\cite{Oroszlany2019,SorianoPalacios2014,He2021TB2J}.

Martinez-Carracedo \emph{et al.} developed a fully relativistic LKAG mapping for nonorthogonal pseudoatomic orbitals, proposed a Hermitian local projection, and demonstrated that a faithful DFT-to-spin mapping rotates the exchange field rather than the complete spinor Hamiltonian~\cite{MartinezCarracedo2023}. Independently, Mankovsky \emph{et al.} treated spin tilts and atomic displacements as perturbations on equal footing in relativistic multiple scattering. Their site-diagonal spin-lattice coefficient contains a torque vertex, a displacement vertex, and two scattering propagators, and it produces a displacement-induced effective magnetic field and torque~\cite{Mankovsky2023SLC}.

The goal here is not to rederive either framework in its native basis. Instead, we formulate their common perturbative structure in an orthonormal spinor Wannier representation and couple it to a screened DFPT perturbation. Wannier interpolation provides a natural basis for both the Hamiltonian and electron-phonon matrix elements~\cite{Marzari2012,Giustino2007,Pizzi2020}, while DFPT supplies the screened first-order lattice response~\cite{Baroni2001}. The electronic calculation outputs a transverse-spin/displacement Hessian. Phonon and magnon spectra are assumed to have been computed independently, and the electronic kernel is projected onto their eigenvectors.

The intended application is a linear one-magnon--one-phonon hybridization in a spin-orbit-coupled magnet. Atomistic spin-lattice Hamiltonians and coupled spin-lattice dynamics provide the broader modeling context. The present linear torque channel differs from a scalar exchange-striction derivative around a collinear state, which commonly produces two-magnon--one-phonon interactions. A total electronic torque response can contain local anisotropy, DMI-like, and symmetric anisotropic intersite contributions without resolving them separately. The direct mode projection therefore avoids an unnecessary spin-model decomposition before constructing the magnon-polaron Hamiltonian.
